\begin{notechapter}
\chapter{Normalization Layers}\label{ch:normalization}
\chaptersubtitle{Batch, layer, instance and group normalization: what they compute and why they help}

\begin{chapterintro}
A normalization layer standardizes the pre-activations of a layer to mean $0$ and variance
$1$, then applies a learned scale and shift. Batch normalization does it with statistics
computed over the mini-batch; layer, instance and group normalization use statistics of a
single example. These layers made very deep networks much easier to train. The reasons are
subtler than the original explanation, and this chapter separates what is proved from what
is conjectured. You will learn
\begin{itemize}[nosep]
  \item the forward pass of batch normalization at training and at test time, and how to fold it
        into the preceding layer;
  \item its backward pass, and why the gradient is a projection;
  \item the scale invariance it creates, and the effective learning rate that results;
  \item the internal covariate shift explanation and the evidence against it;
  \item layer, instance, group and RMS normalization, and which axes each one averages over.
\end{itemize}
\prerequisites{\cref{ch:backpropagation} (the backward pass); \cref{ch:optimizers} for the
discussion of learning rates.}
\end{chapterintro}

\section{Batch normalization}\label{normalization:sec:bn}

Batch normalization (Ioffe and Szegedy, \emph{Batch Normalization: Accelerating Deep Network
Training by Reducing Internal Covariate Shift}, ICML, 2015) acts on each unit separately.
Fix one unit, and let $z_1,\dots,z_m$ be its pre-activations on the $m$ examples of a mini-batch.

\begin{definition}[Batch normalization, training mode]\label{normalization:def:bn}
\[
  \mu_{\mathcal B} = \frac1m\sum_{i=1}^mz_i,\qquad\sigma_{\mathcal B}^2 = \frac1m\sum_{i=1}^m(z_i - \mu_{\mathcal B})^2,\qquad
  \hat z_i = \frac{z_i - \mu_{\mathcal B}}{\sqrt{\sigma_{\mathcal B}^2 + \epsilon}},\qquad y_i = \gamma\hat z_i + \beta,
\]
with a small $\epsilon>0$ and learned parameters $\gamma,\beta$ (one pair per unit).
\end{definition}

For a layer of $n$ units, this is done for each of the $n$ coordinates; for a convolutional
layer, statistics are taken over the batch and all spatial positions of each channel
(\cref{normalization:sec:others}). The normalized output is then passed to the activation.
A bias in the preceding layer is useless, since the mean subtraction removes it; $\beta$ plays its
role.

\paragraph{Why $\gamma$ and $\beta$.} Without them every unit would be forced to have mean $0$ and variance $1$
over the batch, which may not be the best input for the activation (for a sigmoid, it
would keep the unit in its almost linear regime). With $\gamma = \sqrt{\sigma_{\mathcal B}^2 + \epsilon}$ and $\beta = \mu_{\mathcal B}$ the layer
reproduces its input, so normalization does not reduce what the network can represent;
it changes the parametrization, and therefore the optimization. The scale and shift do not
add nonlinearity: for fixed batch statistics, $y_i$ is an affine function of $z_i$.

\paragraph{A nonlinearity in the batch.} During training, $y_i$ depends on the whole batch
through $\mu_{\mathcal B}$ and $\sigma_{\mathcal B}$: the output for one example depends on which other examples share its
batch. The noise this creates acts as a regularizer, and also makes batch normalization
fail for very small batches (with $m = 1$, $\hat z = 0$ and the output is $\beta$ whatever the input).

\begin{example}[Four values]\label{normalization:ex:bn}
Let $z = (1,2,3,6)$, $\gamma = 2$, $\beta = 1$, and $\epsilon = 0$ for readability. Then $\mu_{\mathcal B} = 3$,
$\sigma_{\mathcal B}^2 = \frac14(4 + 1 + 0 + 9) = 3.5$, $\sigma_{\mathcal B} = 1.8708$, and
\[
  \hat z = (-1.0690,\ -0.5345,\ 0,\ 1.6036),\qquad y = 2\hat z + 1 = (-1.1381,\ -0.0690,\ 1,\ 4.2071) .
\]
$\hat z$ has mean $0$ and $\frac14\sum\hat z_i^2 = 1$.
\end{example}

\paragraph{Test time.} At inference, predictions must not depend on the other examples of
the batch, and batches may have size $1$. During training, running estimates of the mean and
variance are kept,
\[
  \mu_{\mathrm{run}}\gets(1 - \alpha)\mu_{\mathrm{run}} + \alpha\mu_{\mathcal B},\qquad\sigma^2_{\mathrm{run}}\gets(1 - \alpha)\sigma^2_{\mathrm{run}} + \alpha\tfrac m{m - 1}\sigma^2_{\mathcal B},
\]
with $\alpha = 0.1$ in PyTorch (its argument \texttt{momentum}; note the convention, where $\alpha$ is the
weight of the \emph{new} value) and the unbiased batch variance. At test time the layer computes
$y = \gamma(z - \mu_{\mathrm{run}})/\sqrt{\sigma^2_{\mathrm{run}} + \epsilon} + \beta$: an affine function of $z$ for each unit, the same for every
example. Since the preceding layer is affine too, the two can be merged
(\cref{normalization:exr:2}): $z = Wa + b$ followed by inference-mode batch normalization equals
$W'a + b'$ with
\[
  W' = \diag\Bigl(\frac\gamma{\sqrt{\sigma^2_{\mathrm{run}} + \epsilon}}\Bigr)W,\qquad b' = \frac{\gamma(b - \mu_{\mathrm{run}})}{\sqrt{\sigma^2_{\mathrm{run}} + \epsilon}} + \beta,
\]
coordinatewise. Deployed networks are often ``folded'' this way. The discrepancy between
batch statistics (training) and running statistics (test) is a classical source of bugs:
forgetting \texttt{model.eval()}, or running statistics corrupted by unusual batches.

\section{The backward pass}\label{normalization:sec:backward}

Every $\hat z_i$ depends on every $z_j$, through $\mu_{\mathcal B}$ and $\sigma_{\mathcal B}$. Let $s = \sqrt{\sigma_{\mathcal B}^2 + \epsilon}$. From $\partial\mu_{\mathcal B}/\partial z_j = 1/m$ and
$\partial\sigma_{\mathcal B}^2/\partial z_j = 2(z_j - \mu_{\mathcal B})/m$, we get $\partial s/\partial z_j = (z_j - \mu_{\mathcal B})/(ms) = \hat z_j/m$, and
\[
  \frac{\partial\hat z_i}{\partial z_j} = \frac1s\Bigl(\delta_{ij} - \frac1m\Bigr) - \frac{z_i - \mu_{\mathcal B}}{s^2}\,\frac{\hat z_j}m
  = \frac1s\Bigl(\delta_{ij} - \frac1m - \frac{\hat z_i\hat z_j}m\Bigr) .
\]

\begin{proposition}[Backward pass of batch normalization]\label{normalization:prop:backward}
Let $g_i = \partial\ell/\partial y_i$. Then $\partial\ell/\partial\beta = \sum_ig_i$, $\partial\ell/\partial\gamma = \sum_ig_i\hat z_i$, and
\[
  \frac{\partial\ell}{\partial z_j} = \frac\gamma{ms}\Bigl(mg_j - \sum_ig_i - \hat z_j\sum_ig_i\hat z_i\Bigr) .
\]
In vector form, $\nabla_z\ell = \frac\gamma s\,P\,g$ with $P = I - \frac1m\mathbf 1\mathbf 1^\top - \frac1m\hat z\hat z^\top$. For $\epsilon = 0$, $P$ is the orthogonal
projection onto the complement of $\operatorname{span}(\mathbf 1, \hat z)$; in particular $\sum_j\partial\ell/\partial z_j = 0$ and
$\sum_j\hat z_j\,\partial\ell/\partial z_j = 0$.
\end{proposition}

\begin{proof}
$\partial\ell/\partial\hat z_i = \gamma g_i$, and $\partial\ell/\partial z_j = \sum_i\gamma g_i\,\partial\hat z_i/\partial z_j$; substituting the Jacobian gives the formula, and
the matrix of the Jacobian is $P/s$, which is symmetric. For $\epsilon = 0$, $\sum_i\hat z_i = 0$ and $\sum_i\hat z_i^2 = m$, so
$\mathbf 1/\sqrt m$ and $\hat z/\sqrt m$ are orthonormal and $P$ is $I$ minus the projection onto their span.
\end{proof}

The projection has a clear meaning: shifting all the $z_i$ by a constant, or scaling them by a
positive factor, leaves $\hat z$ unchanged, so the loss cannot change in the directions $\mathbf 1$ and
$z - \mu_{\mathcal B}\mathbf 1$, and the gradient has no component along them.

\begin{example}[Backward on four values]\label{normalization:ex:backward}
In \cref{normalization:ex:bn}, let the loss be $\ell = y_1 - y_4$, so $g = (1, 0, 0, -1)$. Then $\sum_ig_i = 0$,
$\sum_ig_i\hat z_i = -1.0690 - 1.6036 = -2.6726$ (this is $\partial\ell/\partial\gamma$), and
\[
  \frac{\partial\ell}{\partial z} = \frac2{4\cdot1.8708}\bigl(4g + 2.6726\,\hat z\bigr) = (0.3054,\ -0.3818,\ 0,\ 0.0764),
\]
which sums to $0$ and is orthogonal to $\hat z$; finite differences agree. Treating $\mu_{\mathcal B}$ and $\sigma_{\mathcal B}$ as
constants would give $\gamma g/s = (1.069, 0, 0, -1.069)$, which is wrong. For $z_4$, the direct effect $-1.069$ is almost cancelled by the effect through $\sigma_{\mathcal B}$:
increasing $z_4 = 6$ widens the batch, which shrinks every $\hat z_i$ towards $0$, decreasing $y_4$ and increasing
$y_1$ (both increase $\ell$). The effect through $\mu_{\mathcal B}$ shifts $y_1$ and $y_4$ equally and cancels in $y_1 - y_4$.
\end{example}

\section{Why does it help?}\label{normalization:sec:why}

\paragraph{The original argument.} Ioffe and Szegedy argued that training is slowed by
\emph{internal covariate shift}: as earlier layers are updated, the distribution of the inputs of
each layer changes, and the layer has to keep adapting. Normalizing fixes the first two
moments of each unit's input. Batch normalization does allow much larger learning rates,
reduces sensitivity to the initialization, and keeps activations out of saturated regions.

\paragraph{The evidence.} Santurkar, Tsipras, Ilyas and Madry (\emph{How Does Batch
Normalization Help Optimization?}, NeurIPS, 2018) trained networks where random noise with
changing mean and variance was injected \emph{after} each normalization layer, recreating
covariate shift: they trained almost as well as with plain batch normalization. They also
measured, and proved in simple settings, that batch normalization makes the loss and its
gradient more Lipschitz along the training trajectory, so that larger steps remain
predictive. Their conclusion is that the smoother loss landscape, not reduced covariate
shift, is the main mechanism. The question is not fully settled.

\paragraph{Scale invariance.} A clean, provable effect concerns the weights feeding a
normalized unit. If $z = w^\top a$ is followed by batch normalization (with $\epsilon$ negligible), replacing $w$ by
$cw$ with $c>0$ multiplies all $z_i$ by $c$ and leaves $\hat z$ unchanged: $\ell(cw) = \ell(w)$.

\begin{proposition}[Consequences of scale invariance]\label{normalization:prop:scale}
If $\ell(cw) = \ell(w)$ for all $c>0$, then $\nabla\ell(cw) = \frac1c\nabla\ell(w)$ and $w^\top\nabla\ell(w) = 0$. Gradient descent on $w$ then
increases $\norm w$:
\[
  \norm{w_{k+1}}^2 = \norm{w_k}^2 + \eta^2\norm{\nabla\ell(w_k)}^2 .
\]
\end{proposition}

\begin{proof}
Differentiating $\ell(cw) = \ell(w)$ with respect to $w$ gives $c\nabla\ell(cw) = \nabla\ell(w)$; differentiating with respect to $c$ at
$c = 1$ gives $w^\top\nabla\ell(w) = 0$. Then $\norm{w - \eta\nabla\ell(w)}^2 = \norm w^2 - 2\eta w^\top\nabla\ell(w) + \eta^2\norm{\nabla\ell(w)}^2$.
\end{proof}

Only the direction $u = w/\norm w$ matters, and a step of size $\eta$ on $w$ changes it by about
$\eta\norm{\nabla\ell(w)}/\norm w = \eta\norm{\nabla\ell(u)}/\norm w^2$: the \emph{effective learning rate} for the direction is $\eta/\norm w^2$. It
decreases automatically as $\norm w$ grows, a built-in schedule; weight decay, which shrinks $\norm w$,
then acts as an increase of the effective learning rate rather than as a regularizer (van
Laarhoven, \emph{L2 Regularization versus Batch and Weight Normalization}, 2017).

\section{Other normalizations}\label{normalization:sec:others}

For a convolutional layer the activations form a tensor indexed by the example $n$, the
channel $c$, and the spatial position $(h, w)$. The normalizations differ by the set of entries
over which the mean and variance are computed (\cref{normalization:fig:axes}); all then
apply a learned per-channel (or per-feature) scale and shift.

\begin{center}
\begin{tabular}{lll}
\toprule
method & statistics over & one mean per \\
\midrule
batch norm & $n, h, w$ & channel \\
layer norm & $c, h, w$ & example \\
instance norm & $h, w$ & example and channel \\
group norm ($G$ groups) & $h, w$ and the $C/G$ channels of a group & example and group \\
\bottomrule
\end{tabular}
\end{center}

\begin{figure}[ht]
\centering
\begin{tikzpicture}[x=0.42cm, y=0.42cm]
  \foreach \k/\name/\cols/\rows in {0/batch norm/2/{0,1,2,3}, 1/layer norm/{0,1,2,3,4,5}/1, 2/instance norm/2/1, 3/group norm/{0,1,2}/1}{
    \begin{scope}[xshift=\k*3.7cm]
      \foreach \c in \cols \foreach \r in \rows \fill[bookaccent!45] (\c,-\r) rectangle ++(1,-1);
      \draw[step=1, gray] (0,0) grid (6,-4);
      \draw[thick] (0,0) rectangle (6,-4);
      \node[font=\small] at (3,1) {\name};
      \node[font=\scriptsize] at (3,-4.7) {channels $c$};
      \node[font=\scriptsize, rotate=90] at (-0.8,-2) {examples $n$};
    \end{scope}
  }
\end{tikzpicture}
\caption{Each cell stands for all spatial positions $(h,w)$ of one example and one channel.
The shaded cells form the set over which one mean and one variance are computed: one
channel across the batch (batch norm), one example across channels (layer norm), one
example and one channel (instance norm), one example and a group of channels (group
norm, here $G = 2$).}
\label{normalization:fig:axes}
\end{figure}

\begin{itemize}
  \item \textbf{Layer normalization} (Ba, Kiros and Hinton, 2016) normalizes each example over
        its features, $\hat a_j = (a_j - \mu)/\sqrt{\sigma^2 + \epsilon}$ with $\mu,\sigma^2$ the mean and variance of $(a_1,\dots,a_n)$, then
        $y_j = \gamma_j\hat a_j + \beta_j$. It does not depend on the batch, behaves identically at training and test
        time, and applies to sequences token by token: it is the normalization of recurrent
        networks and transformers.
  \item \textbf{Instance normalization} (Ulyanov, Vedaldi and Lempitsky, 2016) normalizes each
        channel of each image over its spatial positions. It removes the contrast and colour
        statistics of each image, which is useful in style transfer and image generation.
  \item \textbf{Group normalization} (Wu and He, \emph{Group Normalization}, ECCV, 2018) splits the
        channels into $G$ groups and normalizes each group of each example. $G = 1$ is layer
        normalization and $G = C$ instance normalization. It works with batches too small for batch
        normalization, as in detection and segmentation.
  \item \textbf{RMS normalization} (Zhang and Sennrich, \emph{Root Mean Square Layer Normalization},
        NeurIPS, 2019) drops the mean subtraction: $y_j = \gamma_ja_j/\sqrt{\frac1n\sum_ka_k^2 + \epsilon}$. It is cheaper and works as
        well in many transformers.
\end{itemize}
All of these are invariant to rescaling the weights of the preceding layer (for layer
normalization, all weights of the layer at once), so \cref{normalization:prop:scale} applies
to them too.

\begin{pitfall}[Batch normalization and dropout]
Placed before batch normalization, dropout (\cref{ch:dropout}) changes the variance of the
activations between training (with dropout) and test (without), while the running statistics
were collected with dropout: the normalization is then miscalibrated at test time (Li,
Chen, Hu and Yang, \emph{Understanding the Disharmony between Dropout and Batch
Normalization by Variance Shift}, CVPR, 2019). Put dropout after the last batch
normalization layer, or use little of it.
\end{pitfall}

\begin{chaptersummary}
\begin{itemize}[nosep]
  \item Batch norm standardizes each unit over the mini-batch, then applies $\gamma\hat z + \beta$; $\gamma,\beta$ restore the
        capacity to represent any mean and scale. At test time it uses running statistics and is an
        affine map, which folds into the preceding layer.
  \item Its Jacobian is $\frac1sP$, $P$ the projection orthogonal to $\mathbf 1$ and $\hat z$: the gradient ignores shifts and
        rescalings of the batch. Ignoring the dependence of $\mu_{\mathcal B},\sigma_{\mathcal B}$ on $z$ gives wrong gradients.
  \item It allows larger learning rates. The internal covariate shift explanation is contested;
        a smoother loss landscape is the main alternative. Scale invariance makes the gradient
        orthogonal to $w$, $\norm w$ grow, and the effective learning rate $\eta/\norm w^2$.
  \item Layer, instance, group and RMS normalization use statistics of one example; layer norm
        is the standard in transformers.
\end{itemize}
\end{chaptersummary}

\section*{Exercises}
\addcontentsline{toc}{section}{Exercises}

\begin{exercise}\label{normalization:exr:1}
Show that batch normalization in training mode (with $\epsilon = 0$) gives the same output for $(z_1,\dots,z_m)$ and for
$(az_1 + b,\dots,az_m + b)$, $a>0$. What happens for $a<0$?
\solutionlink{normalization:exr:1}
\end{exercise}

\begin{exercise}\label{normalization:exr:2}
Derive the formulas for $W'$ and $b'$ that fold inference-mode batch normalization into the
preceding affine layer. Compute them for one unit with $w = (1,-1)$, $b = 0.5$, $\mu_{\mathrm{run}} = 2$,
$\sigma^2_{\mathrm{run}} = 4$, $\gamma = 3$, $\beta = -1$, $\epsilon = 0$.
\solutionlink{normalization:exr:2}
\end{exercise}

\begin{exercise}\label{normalization:exr:3}
Show directly from the formula of \cref{normalization:prop:backward} that $\sum_j\partial\ell/\partial z_j = 0$, and that the
gradient vanishes when $g$ is an affine function of $\hat z$ ($g_i = c_1 + c_2\hat z_i$). Interpret.
\solutionlink{normalization:exr:3}
\end{exercise}

\begin{exercise}\label{normalization:exr:4}
For one example with features $a = (2, 4, 4, 6)$, compute the outputs of layer normalization
($\gamma = 1$, $\beta = 0$, $\epsilon = 0$) and of RMS normalization ($\gamma = 1$).
\solutionlink{normalization:exr:4}
\end{exercise}

\begin{exercise}\label{normalization:exr:5}
A scale-invariant loss is trained by gradient descent from $\norm{w_0} = 1$ with $\eta = 0.1$, and the gradient norm at
the current direction stays equal to $1$, i.e.\ $\norm{\nabla\ell(w_k)} = 1/\norm{w_k}$. Find the recursion for $\norm{w_k}^2$, and show that
the effective learning rate decays like $1/\sqrt k$.
\solutionlink{normalization:exr:5}
\end{exercise}
\end{notechapter}
